\chapter{Power BI Supervision Dashboard}
\label{chap:powerbi}

\section{A control-room view, not a model host}

The dashboard is where the analytical work becomes usable. I designed it around two operator questions: \emph{what is happening now?} and \emph{what should I inspect?} The first page keeps moisture, laboratory age, process state and the main variables together. The second page opens the evidence behind an abnormal event. This separation avoids turning routine monitoring into a dense diagnostic screen.

Power BI does not run either model. It uses DirectQuery against PostgreSQL views populated by the Python service. The semantic model formats values, applies filters and computes display measures; prediction, novelty scoring and diagnostic transformation have already occurred before the row reaches Power BI.

The Power BI project requests five-second automatic page refresh. Actual refresh can be slower because Power BI waits on DirectQuery, local rendering and source availability \cite{powerbiautorefresh}. The full-screen capture workflow used for the figures below displays a more conservative 60-second presentation label, while the underlying event stream remains five seconds. Neither interval is a confirmed plant historian frequency.

\section{Page 1: Operations Overview}

\Cref{fig:pbi-overview} uses the final high-precision layout. The upper cards distinguish predicted moisture, laboratory moisture, validated error, anomaly score, operating status and severity. The main trend preserves the time relationship among prediction, sparse laboratory markers and anomaly points. Critical variables remain visible below, while the right-hand assessment and guidance panels translate the latest event into an operator-readable summary.

\clearpage
\begin{landscape}
\begin{figure}[p]
\centering
\includegraphics[width=.96\linewidth]{powerbi_overview.png}
\caption{Final Operations Overview capture with high-precision moisture values, process trends and advisory context.}
\source{Current SQL-backed full-screen dashboard capture; Author's work.}
\label{fig:pbi-overview}
\end{figure}
\end{landscape}
\clearpage

The example shown is an injected dashboard-demo event, not a canonical TEST result or a live plant state. In the capture, 0.0798\% predicted moisture is shown separately from the last 0.0752\% laboratory result, and the validated absolute error card refers to a true sample timestamp. The ``LIVE DATA'' pill means that the capture queried the current PostgreSQL view; the source behind that view is still the controlled replay. This distinction is stated here because the interface itself is designed for a future live source.

Three display rules protect interpretation:
\begin{itemize}
\item predicted moisture is a continuous model estimate;
\item laboratory moisture is a sparse reference accompanied by sample time or age; and
\item validated error is calculated only where a represented laboratory sample and prediction coexist.
\end{itemize}

Status colour supports scanning but never carries meaning alone. Green, amber, red and grey are paired with text labels. Numeric precision is deliberately greater for moisture and error than for larger process measurements; this prevents a useful small change from disappearing through formatting while avoiding false precision on every variable.

\section{Page 2: Diagnostics and Investigation}

The diagnostic page in \cref{fig:pbi-diagnostics} begins with a novelty timeline and a selected event. Ranked contributors show which direct variables depart most from their reference profile. The likely subsystem, possible causes and recommended checks appear beside the signature rather than behind a technical model score.

\clearpage
\begin{landscape}
\begin{figure}[p]
\centering
\includegraphics[width=.96\linewidth]{powerbi_diagnostics.png}
\caption{Final Diagnostics capture with event selection, ranked contributors and verification guidance.}
\source{Current SQL-backed full-screen dashboard capture; Author's work.}
\label{fig:pbi-diagnostics}
\end{figure}
\end{landscape}
\clearpage

The captured high-severity demo event points to the drying-air and exhaust circuit. Vacuum, air flow and fan speed are all low relative to the reference profile, so the displayed checks focus on fan demand, damper position, air leakage and vacuum measurement. This is a probable diagnosis for the generated scenario, not proof that one component failed.

Consecutive anomalous rows separated by no more than 30 seconds are grouped into one candidate event. Grouping reduces row-level repetition, but it does not turn a statistical novelty into an approved process alarm. Alarm management requires governed priorities, ownership, response procedures, performance monitoring and alignment with the site philosophy \cite{iec62682}. The present event list is an advisory investigation aid only.

\section{Why the demonstration source is separate}

A short jury session cannot rely on the chance that a useful held-out event happens at the right moment. I therefore prepared a two-day demonstration fork with eight disclosed injected episodes. The source preserves normal rows and sparse laboratory values, and the frozen models are not retrained on the injections.

Across that fork, the model flags 0.071\% of non-injected rows and 58.7\% of injected rows; seven of the eight injected episodes receive a model response. These numbers describe presentation coverage, not independent novelty-model validation. The canonical TEST analysis in Chapter~\ref{chap:models} remains the model evidence and includes the feed-surge miss.

\begin{table}[htbp]
\centering\small
\caption{Dashboard-demo audit and correct interpretation}
\label{tab:dashboard-demo}
\begin{tabularx}{\textwidth}{@{}p{0.30\textwidth}p{0.23\textwidth}Y@{}}
\toprule
Item & Result & Meaning \\
\midrule
Source & 34,560 rows / 2 days & Dedicated presentation replay. \\
Injected scenarios & 8 episodes & Disclosed, deterministic demonstrations. \\
Normal-period model flags & 0.0706\% & Behaviour on non-injected demo rows; not a plant false-positive rate. \\
Injected-period model flags & 58.73\% & Visibility during scripted intervals; not plant detection performance. \\
Episode response & 7 of 8 & Seven demo episodes produce at least one model flag. \\
\bottomrule
\end{tabularx}
\source{\code{resources/dashboard\_demo/MAP\_Dryer\_Dashboard\_Demo\_Audit.json}.}
\end{table}

\section{Usability and presentation performance}

The final launcher opens the report in a PowerPoint-like full-screen presentation workflow so the jury sees the canvas rather than the authoring panes. Page navigation remains manual, but the data state refreshes automatically. Shortening the refresh request to two seconds can make switching pages less smooth because DirectQuery requests and visual rendering overlap. For the defence, stable five-second or conservative presentation refresh is more credible than an aggressive interval that produces lag.

Industrial use would require operator testing, accessible colour and text cues, approved terminology, nuisance-event review and a clear stale-data procedure. It would also require gateway and capacity measurements under expected concurrency. The current pages demonstrate the information architecture and the software connection; they do not establish a validated human--machine interface.
